\documentclass[border=8pt]{standalone}
\usepackage[utf8]{inputenc}
\usepackage{nacal}
\usepackage{tikz}
\usepackage{tikz-3dplot}
\usetikzlibrary{calc,angles,positioning,intersections,quotes,decorations.markings,arrows.meta,decorations.pathreplacing}
\usepackage{tkz-euclide}
\def\Datos{green!40!gray}
\def\Ajuste{blue!50!red}
\def\Residuo{black!25!gray}
\def\Regresor{blue!75!white}
\def\Base{orange!80!black}
\def\Sistem{blue!50!black}
\tikzset{vector/.style={-{Stealth[length=3mm, width=2mm]},very thick},
         basevec/.style={-{Stealth[length=2.2mm, width=1.6mm]},thick,\Base},
         guide/.style={dashed,thin,gray}}

\begin{document}
  \tdplotsetmaincoords{62}{12}
  \begin{tikzpicture}[scale=2.4, tdplot_main_coords]
    \coordinate (O)  at (0,0,0);
    \coordinate (S)  at (1.3,0.35,0);     % parte sistemática realizada: en E
    \coordinate (W)  at (-0.5,0.75,0.8);  % ruido centrado realizado
    \coordinate (We) at (-0.5,0.75,0);    % su proyección sobre E
    \coordinate (V)  at (0.8,1.1,0.8);    % v = s + w
    \coordinate (Ve) at (0.8,1.1,0);      % proyección de v sobre E
    % el subespacio E
    \fill[gray!10,draw=gray!50,very thin]
      (-0.85,-0.35,0) -- (1.9,-0.35,0) -- (1.9,1.6,0) -- (-0.85,1.6,0) -- cycle;
    \node[gray!70!black,anchor=east,font=\small] at (1.85,-0.2,0) {$\mathcal E$};
    \draw[gray!60,thin] (O) -- (0,0,1.2) node[anchor=south,font=\small,gray!70!black] {$\mathcal R$};
    % parte sistemática: dentro de E
    \draw[vector,\Sistem] (O) -- (S)
      node[pos=0.6,anchor=north,font=\small,yshift=-0.5mm] {$\boldsymbol s$};
    % ruido centrado: con componente en E y en R
    \draw[vector,\Residuo] (O) -- (W)
      node[pos=0.75,anchor=east,font=\small,inner sep=2pt] {$\boldsymbol u-\boldsymbol{\mathop{\overline u}}$};
    \draw[guide] (W) -- (We);
    \draw[guide] (O) -- (We);
    % suma
    \draw[vector,\Residuo!70,dashed] (S) -- (V);
    \draw[vector,\Datos] (O) -- (V)
      node[pos=0.5,sloped,above,font=\small,inner sep=1.5pt] {$\boldsymbol v=\boldsymbol s+(\boldsymbol u-\boldsymbol{\mathop{\overline u}})$};
    \draw[guide] (V) -- (Ve);
    \draw[guide] (S) -- (Ve);
    \draw[guide] (O) -- (Ve);
    \node[font=\scriptsize,gray!70!black,anchor=north west,inner sep=1pt] at (Ve) {proyección de $\boldsymbol v$ sobre $\mathcal E$};
    \tkzMarkRightAngle[size=.10,gray,very thin,fill=lightgray!70,opacity=0.6](O,Ve,V)
    % lecturas
  \end{tikzpicture}
\end{document}
